\chapter{RESULTADOS DEL DISEÑO}
La línea primaria académica se desarrolla sobre un corredor horizontal de 6.696 km, con 111 apoyos y 110 vanos. La terna AAAC 50 mm$^2$ demanda 20\,274.361 m de tendido, 240 m de puentes/colas y 616 m de reserva; la cantidad de suministro redondeada es 21\,131 m. Los apoyos se distribuyen en 72 de 13 m, 30 de 15 m y 9 de 17 m, con 111 armados y 160 ramales de retenida. El metrado M01 incorpora el despiece por componente y sus fichas individuales.

La demanda máxima del escenario es 171.486 kW y 190.540 kVA. Con el caso de 10 kV y fuente equivalente a 1.05 pu, la corriente de línea es 11.356 A, la pérdida activa es 2.102 kW y la BT de llegada es 371.067 V. El balance requiere 183.326 kW en generación y 212.481 kVA de fuente. Los resultados incorporan ambos transformadores y auxiliares.

La oferta nominal de 200 kW deja 16.674 kW de margen sin reserva simultánea para Limbani. Con reserva de 25 kW, el déficit es 8.326 kW; la operación aislada debe limitar la carga atendida a la potencia efectivamente disponible. Las alternativas de mayor tensión reducen pérdidas, sin resolver por sí solas la limitación de generación.

El esquema mecánico arroja un cociente máximo de flexión tipo de 0.598 e identifica 46 apoyos con carga ascendente. Se incorporan 16 retenciones adicionales al esquema de estudio. La separación mínima de 7.549 m se refiere al perfil DEM empleado, con las condiciones de interpretación definidas en la memoria.

La entrega reúne en un único expediente la memoria, las ecuaciones, los metrados, el presupuesto referencial y los anexos gráficos. Los libros de cálculo desarrollan las cantidades y verificaciones con fórmulas editables y mantienen sus resúmenes en PDF. Los resultados corresponden a las bases académicas explícitas del documento.
